% Adaptador de compatibilidad del propietario íntegro del radión.
%
% El tratado rector ya carga tipografía, geometría, hipervínculos, teoremas,
% tablas, TikZ y tcolorbox.  Este archivo aporta únicamente los nombres y
% entornos exclusivos del propietario, sin importar su preámbulo autónomo ni
% alterar la jerarquía tipográfica general.

\makeatletter
\@ifundefined{color@RadNavy}{\definecolor{RadNavy}{HTML}{102A43}}{}
\@ifundefined{color@RadBlue}{\definecolor{RadBlue}{HTML}{1F5E7A}}{}
\@ifundefined{color@RadTeal}{\definecolor{RadTeal}{HTML}{168C8C}}{}
\@ifundefined{color@RadCopper}{\definecolor{RadCopper}{HTML}{B66A3C}}{}
\@ifundefined{color@RadGold}{\definecolor{RadGold}{HTML}{D4A73A}}{}
\@ifundefined{color@RadInk}{\definecolor{RadInk}{HTML}{1A202C}}{}
\@ifundefined{color@RadMist}{\definecolor{RadMist}{HTML}{EDF4F7}}{}
\@ifundefined{color@RadCream}{\definecolor{RadCream}{HTML}{FAF7F0}}{}
\@ifundefined{color@RadRose}{\definecolor{RadRose}{HTML}{8D3F55}}{}
\@ifundefined{color@RadGreen}{\definecolor{RadGreen}{HTML}{2F7D5D}}{}
\@ifundefined{color@RadGray}{\definecolor{RadGray}{HTML}{66788A}}{}

\providecommand{\TRIT}{\ensuremath{\mathrm{TRIT}}}
\providecommand{\HMT}{\ensuremath{\mathrm{HMT}}}
\providecommand{\Rstate}{\mathcal R_{12}}
\providecommand{\epsR}{\varepsilon_R}
\providecommand{\epsone}{\varepsilon_R^{(1)}}
\providecommand{\pih}{\pi_{\!\HMT}}
\providecommand{\alphah}{\alpha_{\!\HMT}}
\providecommand{\hret}{\hbar_{\mathrm{ret}}}
\providecommand{\hfull}{h_{\mathrm{ret}}}
\providecommand{\diff}{\mathop{}\!\mathrm d}
\providecommand{\e}{\mathrm e}
\providecommand{\R}{\mathbb R}
\providecommand{\C}{\mathbb C}
\providecommand{\F}{\mathbb F}
\providecommand{\Z}{\mathbb Z}
\providecommand{\dd}{\,\mathrm d}
\providecommand{\status}[1]{\textsf{\bfseries #1}}
\providecommand{\sourcepath}[1]{\nolinkurl{#1}}

\@ifundefined{typebox}{%
  \newtcolorbox{typebox}[1][]{enhanced,breakable,colback=white,
    colframe=RadBlue,boxrule=.7pt,arc=1.5mm,left=4mm,right=4mm,
    top=3mm,bottom=3mm,title={Recibo de tipos},fonttitle=\bfseries,
    coltitle=RadNavy,#1}%
}{}
\@ifundefined{narrativebox}{%
  \newtcolorbox{narrativebox}[1][]{enhanced,breakable,colback=white,
    colframe=RadGold,borderline west={2.4pt}{0pt}{RadGold},
    boxrule=.25pt,arc=0mm,left=5mm,right=4mm,top=3mm,bottom=3mm,#1}%
}{}

% Las once portadillas internas del propietario autónomo se convierten en
% divisiones científicas ordinarias dentro del capítulo 52.  Se evita así
% introducir partes autónomas, gigantismo tipográfico o páginas vacías.
\providecommand{\partopener}[3]{\section{#2}\noindent\emph{#3}\par\medskip}

\providecommand{\BigRadionEquation}{%
  \begin{equation*}
  \boxed{\displaystyle
  \frac{180^{\circ}}{\pih}
  =50^{\circ}+A^{\circ}-\frac{20}{81}\,\Delta_4^{\circ}+\epsR^{\circ}}
  \end{equation*}}
\makeatother
